\documentclass[conference]{IEEEtran}

% ── Packages ────────────────────────────────────────────────────────────────
\usepackage[T1]{fontenc}
\usepackage[utf8]{inputenc}
\usepackage{cite}
\usepackage{amsmath,amssymb}
\usepackage{graphicx}
\usepackage{booktabs}
\usepackage{array}
\usepackage{url}
\usepackage{hyperref}
\usepackage{xcolor}
\usepackage{multirow}
\usepackage{caption}

\hypersetup{
    colorlinks = true,
    linkcolor  = blue,
    citecolor  = blue,
    urlcolor   = blue,
}

% ── Figure path ─────────────────────────────────────────────────────────────
\graphicspath{{../QKD_Simulation/Advanced Protocol Set/plots/findings/}}

% ── Title ───────────────────────────────────────────────────────────────────
\title{Simulation-Based Comparative Benchmarking of QKD Protocols:\\
BB84, E91, and MDI-QKD Over Realistic Optical Fibre Channels}

\author{
  \IEEEauthorblockN{[Author Name(s)]}
  \IEEEauthorblockA{
    [Department Name]\\
    [University Name]\\
    [City, Country]\\
    [email@university.edu]
  }
}

\begin{document}

\maketitle

% ════════════════════════════════════════════════════════════════════════════
\begin{abstract}
Quantum Key Distribution (QKD) allows two parties to exchange cryptographic
keys with security guaranteed by the laws of quantum mechanics rather than
computational hardness assumptions. While several QKD protocols have been
proposed and studied theoretically, objective performance comparisons under
identical and realistic channel conditions remain scarce in the open literature.
This paper presents a simulation-based comparative benchmark of three
prominent QKD protocols---BB84, E91, and MDI-QKD---implemented over a
realistic optical fibre channel model that incorporates photon loss,
depolarisation noise, and finite detector efficiency. Eight standardised
performance metrics are evaluated across 30 independent simulation runs at
a reference distance of 10~km, and a distance sweep from 1~km to 50~km
characterises scalability behaviour. Results show that BB84 achieves the
highest throughput (164\,835~bits/s) and lowest communication overhead
(51~messages), while MDI-QKD provides the strongest security guarantees at
the cost of a 12$\times$ increase in classical overhead. E91's
entanglement-based security proof is confirmed via a Bell parameter of
$S \approx 2.60 > 2$, though it suffers the highest channel loss rate
(72.1\%) due to its requirement for coincident photon detection at both
endpoints. These findings offer protocol-selection guidance for quantum
network designers balancing security requirements against operational
efficiency.
\end{abstract}

\begin{IEEEkeywords}
Quantum Key Distribution, BB84, E91, MDI-QKD, QBER, quantum network
simulation, Bell inequality, comparative benchmarking, optical fibre channel
\end{IEEEkeywords}

% ════════════════════════════════════════════════════════════════════════════
\section{Introduction and Literature Review}

The security of classical cryptographic systems such as RSA depends on the
practical difficulty of solving certain mathematical problems. Advances in
quantum computing, however, threaten to render these assumptions insufficient.
In 1994, Shor demonstrated that a sufficiently powerful quantum computer
could factorise large integers in polynomial time \cite{shor1994}, placing
widely deployed public-key infrastructure at risk. Quantum Key Distribution
offers an alternative whose security rests on physical law: an eavesdropper
who intercepts a quantum channel necessarily disturbs the transmitted states,
and this disturbance can be detected by the legitimate parties \cite{gisin2002}.

\subsection{BB84}

The first QKD protocol was introduced by Bennett and Brassard in 1984
\cite{bb84}. Alice encodes each key bit into the polarisation of a single
photon using one of two conjugate bases (rectilinear $\{|0\rangle,|1\rangle\}$
or diagonal $\{|+\rangle,|-\rangle\}$), chosen randomly. Bob measures each
arriving photon in a randomly selected basis. After transmission, Alice and
Bob publicly compare their basis choices and retain only those bits where they
selected the same basis---the \emph{sifted key}. Any eavesdropping
introduces a detectable Quantum Bit Error Rate (QBER). Security proofs for
BB84 against arbitrary attacks were established by Mayers \cite{mayers2001}
and Lo and Chau \cite{lo1999}. BB84 has since been demonstrated experimentally
over distances exceeding 400~km in optical fibre \cite{boaron2018}.

\subsection{E91}

Ekert proposed an entanglement-based QKD protocol in 1991 \cite{e91}.
An entanglement source generates pairs of photons in a Bell state and
distributes one photon to Alice and the other to Bob. Each party independently
measures their photon in a randomly chosen basis from a set of three angles.
The subset of measurements performed in \emph{different} bases are used to
evaluate the CHSH Bell inequality \cite{chsh1969}; a value $S > 2$ confirms
entanglement and proves the absence of any local-hidden-variable eavesdropping
strategy. The subset where Alice and Bob chose the \emph{same} basis forms the
key. E91 is notable because its security proof is device-independent in the
sense that it does not rely on trusting the source of entangled photons.

\subsection{MDI-QKD}

A practical vulnerability in deployed QKD systems is the detector
side-channel attack: imperfect physical detectors can leak information
about Bob's measurement outcomes to an eavesdropper \cite{lydersen2010}.
Lo, Curty, and Qi addressed this by introducing
Measurement-Device-Independent QKD (MDI-QKD) in 2012 \cite{mdi2012}. Both
Alice and Bob send quantum states to an untrusted relay node (Charlie), who
performs a Bell-State Measurement (BSM) and publicly announces the result.
Because neither party trusts Charlie's detectors, all detector side-channel
attacks are eliminated by construction. MDI-QKD has been experimentally
realised over metropolitan-scale networks \cite{tang2016}.

\subsection{Simulation-Based Evaluation}

Physical implementations of QKD require single-photon sources, precision
synchronisation, and expensive cryogenic detectors, making experimental
comparison of multiple protocols impractical for most research groups.
Simulation provides a reproducible and cost-effective alternative. Doolittle
and Cubeddu \cite{doolittle2026} propose a roadmap for quantum network
simulation and emulation, arguing that rigorous software-level modelling of
physical-layer imperfections is essential for realistic protocol evaluation.
NetSquid \cite{coopmans2021} provides a discrete-event quantum network
simulator with hardware-level noise models and has been used for end-to-end
QKD network studies. Despite these tools, no prior work has benchmarked
BB84, E91, and MDI-QKD under \emph{identical} channel conditions using a
unified metric framework---the gap this paper addresses.

% ════════════════════════════════════════════════════════════════════════════
\section{Methodology}

\subsection{Simulation Framework}

The simulation is built in Python 3.12. The \emph{Basic Protocol Set}
implements all three protocols as self-contained modules with no external
quantum-simulator dependency, enabling fast multi-run statistical collection.
The \emph{Advanced Protocol Set} uses NetSquid 1.1.x for discrete-event
simulation with hardware-accurate models. All benchmarking results reported
in this paper are produced by the Basic Protocol Set under the channel
parameters described below.

\subsection{Channel and Device Models}

Each simulation instantiates the same physical channel:

\begin{itemize}
  \item \textbf{Fibre loss}: $0.1$~dB initial loss plus $0.2$~dB/km length-dependent
        attenuation (standard SMF-28 fibre), giving a transmission probability
        $\eta = 10^{-(0.1 + 0.2 \cdot d)/10}$ for distance $d$ km.
  \item \textbf{Depolarisation noise}: random bit-flip applied with rate
        $p_\text{depol} = 0.008$ per qubit traversal.
  \item \textbf{Propagation delay}: speed of light in fibre $c = 1.9 \times 10^5$~km/s.
  \item \textbf{Detector efficiency}: $\eta_d = 0.85$ for all single-photon detectors.
  \item \textbf{E91 entanglement source}: generation efficiency 0.75, pair
        generation success rate 0.85, entanglement fidelity 0.92.
  \item \textbf{MDI-QKD BSM node}: two-photon Bell-State Measurement with
        distance-dependent success rate
        $p_\text{BSM}(d) = 0.25\,(1-0.005d)(1-0.003d)$, capped at a minimum
        of 0.05.
\end{itemize}

\subsection{Protocol Implementations}

\textbf{BB84} follows the standard four-state prepare-and-measure scheme.
Alice prepares 1\,000 qubits per run in randomly selected basis-bit
combinations and transmits them through the quantum channel. Bob measures in
a randomly chosen basis; after basis sifting, a random subset of the sifted
key is used for QBER estimation, followed by error correction and privacy
amplification.

\textbf{E91} distributes 1\,000 entangled pairs per run from a central source
placed equidistant from Alice and Bob. Each party measures in one of three
angles: Alice uses $\{0°, 22.5°, 45°\}$ and Bob uses $\{-22.5°, 0°, 22.5°\}$.
Measurement correlations are computed as $E(\theta_a, \theta_b)=\cos
(2(\theta_a - \theta_b))$ for maximally entangled states, degraded by
the entanglement fidelity parameter. The CHSH parameter is computed as
$S = |E_{00} + E_{02} + E_{22} - E_{20}|$ where subscripts index the angle
pair choices.

\textbf{MDI-QKD} has Alice and Bob each prepare and send 1\,000 qubits per
run toward the central relay Charlie. Charlie attempts a BSM on each arriving
pair; successful measurements with matching bases are announced and used for
key sifting. Error correction and privacy amplification are applied to
the resulting sifted key.

\subsection{Benchmarking Protocol}

To ensure statistical validity, each protocol is executed $N=30$ independent
times at the reference distance of 10~km. Means and standard deviations are
computed per metric. Additionally, a distance sweep over $\{1, 5, 10, 20, 30,
50\}$~km is conducted with $N=8$ runs per distance point to characterise
scalability and metric degradation trends.

\subsection{Performance Metrics}

Eight metrics are recorded after each run:

\begin{enumerate}
  \item \textbf{Raw Key Rate}: sifted key length in bits after basis reconciliation.
  \item \textbf{QBER}: fraction of key bits in error ($\times 100\%$).
  \item \textbf{Latency}: wall-clock simulation time in milliseconds.
  \item \textbf{Channel Loss Rate}: fraction of transmitted qubits not received.
  \item \textbf{Throughput}: final secure key bits per second.
  \item \textbf{Communication Overhead}: number of classical messages exchanged.
  \item \textbf{Synchronisation Time}: time to align measurement bases in milliseconds.
  \item \textbf{Computation Time per Round}: per-qubit processing time in nanoseconds.
\end{enumerate}

% ════════════════════════════════════════════════════════════════════════════
\section{Results}

\subsection{Metric Distributions at 10 km}

Table~\ref{tab:results} summarises all eight metrics across 30 simulation
runs at 10~km. Fig.~\ref{fig:boxplots} shows the full statistical
distributions as box plots.

\begin{table*}[t]
\caption{Performance Metric Summary: Mean $\pm$ Standard Deviation (30 runs, 10 km fibre)}
\label{tab:results}
\centering
\renewcommand{\arraystretch}{1.25}
\begin{tabular}{l c c c}
\toprule
\textbf{Metric} & \textbf{BB84} & \textbf{E91} & \textbf{MDI-QKD} \\
\midrule
Raw Key Rate (bits)           & $170.2 \pm 11.2$ & $29.6 \pm 3.7$   & $19.7 \pm 2.3$    \\
QBER (\%)                     & $0.00 \pm 0.00$  & $0.00 \pm 0.00$  & $4.33 \pm 1.97$   \\
Latency (ms)                  & $1.03 \pm 0.02$  & $1.05 \pm 0.05$  & $1.43 \pm 0.13$   \\
Channel Loss Rate (\%)        & $38.5 \pm 1.3$   & $72.1 \pm 1.5$   & $39.7 \pm 1.3$    \\
Throughput (bits/s)           & $164\,835 \pm 10\,324$ & $28\,373 \pm 3\,495$ & $13\,850 \pm 1\,751$ \\
Communication Overhead (msgs) & $51 \pm 0$       & $131 \pm 0$      & $593 \pm 12$      \\
Synchronisation Time (ms)     & $1.26 \pm 0.00$  & $1.52 \pm 0.00$  & $2.85 \pm 0.00$   \\
Computation Time/Round (ns)   & $14.8 \pm 5.1$   & $387.9 \pm 25.4$ & $67.3 \pm 21.7$   \\
\bottomrule
\end{tabular}
\end{table*}

\begin{figure}[h]
  \centering
  \includegraphics[width=\columnwidth]{Fig1_BoxPlots.png}
  \caption{Box-plot distributions of all eight performance metrics across
           30 runs per protocol at 10~km. Horizontal lines inside boxes
           denote medians; boxes span the IQR; whiskers extend to 1.5$\times$IQR.}
  \label{fig:boxplots}
\end{figure}

\subsection{Key Rate and Throughput}

BB84 achieves a mean raw key rate of 170.2~bits per run---5.7$\times$ higher
than E91 (29.6~bits) and 8.6$\times$ higher than MDI-QKD (19.7~bits). This
translates directly into throughput: BB84 at 164\,835~bits/s versus
13\,850~bits/s for MDI-QKD, a 12$\times$ gap. The primary cause of E91's
lower key rate is its high channel loss: because both photons of an entangled
pair must survive their respective half-channels and be detected
simultaneously, the effective transmission probability is $\eta^2$ rather
than $\eta$. At 10~km with $\eta \approx 0.53$, this gives a coincidence
probability of approximately 0.28, which combined with the 85\% detector
efficiency explains the 72.1\% channel loss rate observed.

MDI-QKD's lower key rate stems from the BSM success rate, which is
distance-dependent and averages around 10--11\% at 10~km in this simulation.
Of the 1\,000 qubits each party sends, only those where Charlie's BSM
succeeds \emph{and} Alice and Bob's bases match contribute to the sifted key.

\subsection{Quantum Bit Error Rate}

BB84 and E91 both achieve 0.00\% QBER at 10~km in this simulation. This
reflects the noise model parameters: at 10~km, the depolarisation rate and
detector error combine to a level below the threshold that triggers errors
in the sifted key under these parameter choices. MDI-QKD records a
mean QBER of 4.33\% $(\pm1.97\%)$, attributable to the interference
visibility limitations of the BSM node (95\% in this model) and the
distance-dependent measurement error rate of Charlie's detectors. The
MDI-QKD QBER remains below the 11\% security threshold at which privacy
amplification becomes infeasible under standard BB84-type security analysis.

\subsection{Communication Overhead}

The most striking operational difference between the three protocols is
classical communication overhead. BB84 requires 51 messages per run, covering
session establishment, basis reconciliation, error correction, and key
confirmation. E91 requires 131 messages---the additional cost of three-party
entanglement synchronisation, Bell-test coordination, and the more complex
post-processing that the entanglement-based security proof requires.
MDI-QKD demands 593 messages per run, a 12$\times$ increase over BB84.
This stems from Charlie's BSM announcements (one per successful measurement),
additional three-party synchronisation, and the extra authentication rounds
required to keep the relay node honest.

\subsection{Security-Performance Trade-off}

Fig.~\ref{fig:tradeoff} illustrates the core trade-off. Each point represents
one simulation run; diamond markers denote per-protocol means. BB84 clusters
at high throughput with near-zero QBER. E91 occupies a mid-throughput
position with zero QBER, while MDI-QKD sits at the lowest throughput but
with the strongest practical security guarantee (immune to all detector
side-channel attacks regardless of QBER).

\begin{figure}[h]
  \centering
  \includegraphics[width=\columnwidth]{Fig5_Security_Performance.png}
  \caption{Security--performance scatter plot: QBER (lower = more secure, x-axis)
           versus throughput (higher = faster, y-axis). Each point is one run;
           diamonds are per-protocol means.}
  \label{fig:tradeoff}
\end{figure}

\subsection{E91 Bell Parameter}

The CHSH Bell parameter computed across runs averages $S \approx 2.60$,
exceeding the local-realism bound of $S = 2$ and the theoretical
maximum for classical correlations. This confirms that the entangled states
distributed in the simulation retain sufficient fidelity to provide
a quantum-mechanical security guarantee. The quantum mechanical maximum
is $S = 2\sqrt{2} \approx 2.83$; the gap between our result and this
maximum reflects the 92\% entanglement fidelity of the source model and
the depolarisation noise applied during transmission.

\subsection{Distance Analysis}

Fig.~\ref{fig:distance} shows how the eight metrics evolve over 1--50~km.
All protocols degrade with distance, but at different rates.

\begin{figure}[h]
  \centering
  \includegraphics[width=\columnwidth]{Fig2_Distance_Analysis.png}
  \caption{Distance analysis: each metric averaged over 8 runs per distance
           point. Shaded bands denote $\pm 1$ standard deviation. All protocols
           use identical channel models.}
  \label{fig:distance}
\end{figure}

Raw key rate decreases monotonically for all three protocols. BB84 degrades
most gracefully, retaining positive key generation beyond 50~km in most runs.
E91's coincidence-detection requirement causes its key rate to approach zero
faster; beyond 30~km, many runs produce insufficient coincidences for key
generation. MDI-QKD's distance degradation is intermediate: the
mid-point relay halves the effective distance each photon travels, but
the BSM success rate penalty offsets this advantage at longer distances.

\subsection{Normalised Comparison}

Fig.~\ref{fig:radar} presents a radar chart in which all metrics are
normalised to [0,\,1] and oriented so that \emph{outward = better}.
Metrics where lower is preferred (QBER, loss rate, overhead, latency) are
inverted before plotting. BB84 dominates the efficiency axes; E91 leads on
latency and QBER; MDI-QKD trails on most dimensions but represents the only
protocol that eliminates detector side-channel attacks without any
trust assumptions on measurement devices.

\begin{figure}[h]
  \centering
  \includegraphics[width=0.85\columnwidth]{Fig3_Radar_Chart.png}
  \caption{Normalised performance radar chart. Each axis is scaled to [0,\,1]
           with outward = better. Axes for QBER, channel loss rate,
           communication overhead, latency, synchronisation time, and
           computation time are inverted.}
  \label{fig:radar}
\end{figure}

% ════════════════════════════════════════════════════════════════════════════
\section{Discussion}

The results allow three protocol-selection guidelines to be stated:

\textbf{(1) High-bandwidth, trusted-device environments.} BB84 is the
appropriate choice. Its simplicity translates into the highest key rate,
lowest overhead, and fastest computation per round. For networks where
detectors are physically secured, BB84's efficiency advantage outweighs
the absence of detector-independence guarantees.

\textbf{(2) Scenarios requiring a quantum-mechanical security proof.}
E91 is preferred when the security proof must rest on a verifiable physical
principle (Bell inequality violation) rather than assumptions about the
implementation of Alice's and Bob's devices. The 5.7$\times$ key-rate
penalty relative to BB84 and the high channel loss rate are the costs of
this stronger guarantee.

\textbf{(3) Untrusted or adversarially controlled measurement devices.}
MDI-QKD is the only option when detector side-channel attacks are a
credible threat---for example, when QKD equipment is supplied by a
potentially adversarial vendor or operated in a shared facility. The
12$\times$ communication overhead and 8.6$\times$ key-rate penalty are
acceptable trade-offs for the complete elimination of detector attacks.

A limitation of this study is that the simulation models photon loss and
depolarisation noise but does not include phase noise, timing jitter, or
multi-photon pulse attacks, which affect real optical implementations.
Future work will incorporate these effects and extend the comparison to
twin-field QKD \cite{lucamarini2018}, which achieves longer distances by
relaxing the photon-arrival timing requirement at the relay node.

% ════════════════════════════════════════════════════════════════════════════
\section{Conclusion}

This paper presented a simulation-based comparative benchmark of three
QKD protocols---BB84, E91, and MDI-QKD---under identical realistic
optical fibre channel conditions. Thirty independent simulation runs at
10~km, combined with a six-point distance sweep from 1~km to 50~km,
provide statistically robust evidence for the following conclusions:

\begin{itemize}
  \item BB84 achieves the highest efficiency: 170.2~bits raw key rate,
        164\,835~bits/s throughput, and only 51 classical messages per session.
  \item E91's Bell parameter of $S \approx 2.60$ confirms entanglement-based
        security, but coincident-detection requirements yield a 72.1\% channel
        loss rate at 10~km.
  \item MDI-QKD eliminates detector side-channel attacks at the cost of
        593 classical messages per session and a 12$\times$ throughput
        reduction relative to BB84.
  \item All protocols degrade predictably with distance; MDI-QKD degrades
        fastest beyond 20~km due to compounding photon-loss and BSM
        success-rate penalties.
\end{itemize}

These findings provide concrete, reproducible guidance for quantum network
designers selecting between the three dominant QKD paradigms. The simulation
code and all benchmark scripts are publicly available in the project
repository.

% ════════════════════════════════════════════════════════════════════════════
\bibliographystyle{IEEEtran}
\bibliography{references}

\end{document}
